% Source-derived chapter 06: Mirror of  upward flows
% Original source: very-stable-higgs-bundles-equivariant-multiplicity-mirror-symmetry
\chapter{Mirror of  upward flows}
\label{very-stable-higgs-bundles-equivariant-multiplicity-mirror-symmetry-chapter-06}
\source{very-stable-higgs-bundles-equivariant-multiplicity-mirror-symmetry}

\begin{remark}[Source section]
\label{very-stable-higgs-bundles-equivariant-multiplicity-mirror-symmetry-chapter-06-source}
\source{very-stable-higgs-bundles-equivariant-multiplicity-mirror-symmetry}
\source{very-stable-higgs-bundles-equivariant-multiplicity-mirror-symmetry}
This chapter reproduces the corresponding section of the retrieved
LaTeX source, including its definitions, statements, examples, and
proofs. It has not yet been translated into Lean.
\notready
\end{remark}

% The source section title was: Mirror of  upward flows
\label{mirror}
\source{very-stable-higgs-bundles-equivariant-multiplicity-mirror-symmetry}

In this section we consider the mirror of $\calO_{W^+_\calE}$ the structure sheaf of a very stable type $(1,...,1)$ upward flow. We recall that we work with $\GL_n$-Higgs bundles. Because $\GL_n$ is its own Langlands dual we expect the mirror of $\calO_{W^+_\calE}$ to be a coherent sheaf on $\M$, or a complex of sheaves which -- generically at least -- should be a fibrewise Fourier-Mukai transform on the self-dual Jacobian of a smooth spectral curve. We start by recalling these notions. 
\subsection{Jacobians, duality, Fourier-Mukai transform}
The Jacobian $J:=J_0(C)$ of a smooth complex projective curve $C$ is the moduli space of degree $0$ line bundles on $C$. It is an Abelian variety with multiplication of line bundles. %There are universal line bundles on $\L\to J\times C$ well-defined up to tensorisation with a line bundle on $J$. Sometimes it is convenient to normalize the universal bundle by picking a point $c\in C$, and insisting that $\L|_{\{c\}}\cong \calO_J$ is trivial
Its dual Abelian variety is defined as $\Pic^0(J)$ the moduli space of translation invariant line bundles on $J$. One can identify $\Pic^0(J)$ with $J$ in two different ways yielding the same isomorphism between the two. 

First, we fix a point $c_0\in C$ and define the {\em Abel-Jacobi map} as $\alpha_{c_0}:C\to J$ by $$\alpha_{c_0}(c):=\calO(c-c_0).$$ Then it turns out \cite[Lemma 11.3.1]{birkenhake-lange} that the  restriction map \beq\label{ajiso} \alpha_{c_0}^*:\Pic^0(J)\to J\eeq is an isomorphism. 
\source{very-stable-higgs-bundles-equivariant-multiplicity-mirror-symmetry}

%More generally, if $M\in J_{-1}(C)$ a degree $-1$ line bundle we can define \bes \begin{array}{cccc} a_M:&C&\to& J \\ & c &\mapsto & M(c).\end{array} \ees In particular, $a_{\calO(c_0)}=a_{c_0}$. 

%If $M_1\in J_{-1}(C)$ is another degree $-1$ line bundle on $C$ then $$\alpha_{M_1}= t_{MM_1^{-1}}\circ \alpha_{M},$$
%where $t_{MM_1^{-1}}:J\to J$ is $L\mapsto LMM_1^{-1}$ the translation by $MM_1^{-1}$. Thus $$\alpha_{M_1}^*(L)=\alpha_{M}^*(t_{MM_1^{-1}}^*(L))$$ as $L\in \Pic^0(J)$ are precisely the translation invariant line bundles on $J$, we get $\alpha_{M}^*(L)=\alpha_{M_1}^*(L)$ thus the isomorphism \eqref{ajiso} is independent of $c\in C$, and indeed agrees with $a_M^*$ for all $M\in J_{-1}$.  

Another isomorphism between $\Pic^0(J)$ and $J$ is given by $\calL$ the line bundle of the Theta divisor $\Theta$.  One defines \beq \label{philiso} \phi_\calL: J\to \Pic^0(J)\eeq by $L\mapsto t_L^*(\calL) \otimes \calL^*$, where $t_L:J\to J$ is translation by $L$. The map \eqref{philiso} is an isomorphism as $\calL$ defines  a principal polarization on $J$. It follows from \cite[Proposition 11.3.5]{birkenhake-lange} that $-\phi_\calL=(\alpha^*_{c_0})^{-1}$. Thus $-\phi_\calL$ and $\alpha^*_{c_0}$ defines the same isomorphism between $J$ and $\Pic^0(J)$. In particular, $\alpha^*_{c_0}$ is independent of the choice of $c_0\in C$.  We will identify $J$ with $\Pic^0(J)$ using this isomorphism below. 
\source{very-stable-higgs-bundles-equivariant-multiplicity-mirror-symmetry}

We have the {\em Poincar\'e bundle} $\bP$ on $J\times \Pic^0(J)$ unique with the properties that \beq \label{poiniden} \bP|_{J\times \{L\}}\cong L\eeq for every $L\in \Pic^0(J)$ and $$\bP|_{\{\calO_C\}\times \Pic^0(J)}\cong \calO_{\Pic^0(J)}.$$
\source{very-stable-higgs-bundles-equivariant-multiplicity-mirror-symmetry}
For $L\in J$ we will denote the line bundle $$\bP_{L}:= \bP|_{\{L\}\times \Pic^0(J)}$$ on $\Pic^0(J)\cong J$, where we identified  $\Pic^0(J)$ and $J$ with $\alpha_{c_0}^*$. This translation invariant line bundle on $\Pic^0(J)$ is characterised by the property that \beq\label{char}\alpha^*_{c_0}(\bP_L)\cong L.\eeq
\source{very-stable-higgs-bundles-equivariant-multiplicity-mirror-symmetry}
In particular, for ${L_1},L_2\in J$ we have \beq \label{tensorp}\bP_{{L_1}L_2}\cong  \bP_{{L_1}}\otimes \bP_{L_2},\eeq and $$\bP_{\calO_C}=\calO_{J}.$$
\source{very-stable-higgs-bundles-equivariant-multiplicity-mirror-symmetry}
Using the Poincar\'e bundle $\bP$ on $J\times \Pic^0(J)\cong J\times J$  we can define the Fourier-Mukai transform \cite{mukai}. For a bounded complex $\calF^\bullet \in D_{coh}^b(J)$ of coherent sheaves on $J$ we denote by $$S(\calF^\bullet):=(\pi_2)_*(\pi_1^*(\calF^\bullet)\otimes^L \bP)\in D^b_{coh}(J),$$ where $\pi_1$ and $\pi_2$ are the projections of $J\times J$ to the two factors. Then $S$ is called
the {\em Fourier-Mukai transform }of $\calF^\bullet$. For example we have $$S(\calO_{L})=\bP_L.$$ 

We recall three properties of the Fourier-Mukai transform. First by \cite[Theorem 3.13.1]{mukai} for any $\calF^\bullet \in D^b_{coh}(J)$ \beq \label{skewinv} S(S(\calF^\bullet))=(-1)^*(\calF^\bullet )[-g], \eeq where $-1:J\to J$ denotes the inverse map on the abelian variety $J$ and $g=\dim(J)$ is the genus of $C$.   The  second \cite[Theorem 3.13.4]{mukai} \beq \label{skewdual}  S(\calF^{\bullet\vee}) \cong (-1)^*S(\calF^\bullet)^\vee [g].\eeq Here $(\calF^\bullet)^\vee:=R{\mathcal Hom}(\calF^\bullet,\calO_J)$ is the derived dual complex. Finally, by proper base change and \eqref{poiniden} we see that the  restriction of the Fourier-Mukai transform to the identity satisfies \beq \label{fmiden} S(\calF^\bullet)|_{\{\calO_C\}}\cong R\pi_*(\calF^\bullet), \eeq
\source{very-stable-higgs-bundles-equivariant-multiplicity-mirror-symmetry}
where $\pi:J\to \{\calO_C\}$ is the projection to the point. 

We can pull back \beq \label{univlambda}\L:=(\alpha_{c_0}\times \alpha_{c_0}^{*-1})^*(\bP)\eeq the Poincar\'e bundle on $J\times \Pic_0(J)$ to $C\times J$ via the map $$\alpha_{c_0}\times \alpha_{c_0}^{*-1}:C\times J\to J\times \Pic_0(J)$$ to get a universal line bundle on $C\times J$ satisfying \beq \label{univl}\L|_{C\times \{L\}}\cong L\eeq for all $L\in J$ and normalized
\source{very-stable-higgs-bundles-equivariant-multiplicity-mirror-symmetry}
by an isomorphism \beq\label{normalize} \bbbeta:\L|_{\{c_0\}\times J}\stackrel{\cong}{\to}\calO_{J}.\eeq
\source{very-stable-higgs-bundles-equivariant-multiplicity-mirror-symmetry}

\begin{remark}
\source{very-stable-higgs-bundles-equivariant-multiplicity-mirror-symmetry}
\notready \label{framedjacobian} As a universal line bundle on $C\times J$  the line bundle $\L$ is not unique as we can tensor it with the pull back of any line bundle on $J$, and the universal property still holds. However
\source{very-stable-higgs-bundles-equivariant-multiplicity-mirror-symmetry}
	$(\L,\bbbeta)$ is the unique universal object for the functor assigning to a scheme $T$  families of framed line bundles $(L,\beta)$, where $L$ is a line bundle on $C\times T$ and $\beta:L|_{\{c_0\}\times T}\cong \calO_T$. 

If $c^\prime_0\in C$ is another point, then $\L\otimes \pi_2^*\left(\bP_{\calO(c_0-c^\prime_0)}\right)|_{\{c^\prime_0\}\times J}\cong \calO_J$.  This is because  $\L|_{\{c^\prime_0\}\times J}\cong (\alpha_0^{*-1})^*\bP_{\calO(c^\prime_0-c_0)}$ from \eqref{char}. Therefore the universal bundle normalised at $c^\prime_0$ is given by \beq \label{changenorm} \L^\prime=\L\otimes \pi_2^*\left(\bP_{\calO(c_0-c^\prime_0)}\right). \eeq 
\source{very-stable-higgs-bundles-equivariant-multiplicity-mirror-symmetry}
\end{remark}

%In fact, any line bundle on $C\times J$ satisfying \eqref{univl} is of the form $\L$ for some $M\in J_{-1}$. In particular if $N\in J$ is a degree $0$ line bundle on $C$, then \beq \label{translateL}\L^{NM}\cong \bP_N\otimes \L\eeq

Finally, we denote by $\pi:\mathbf{C}\to C$  a finite morphism of smooth projective curves, by  $\tilde{J}$  the Jacobian of $\mathbf{C}$ and  $\tilde{\bP}$ denotes the Poincar\'e bundle on $\tilde{J}\times \tilde{J}$. Then, in this setting,  we have the norm map ${\mathrm{Nm}}_\pi:\tilde{J}\to J$ by \beq \label{norm}L\mapsto \det(\pi_*(M))\otimes \det(\pi_*(\calO_{\mathbf{C}}))^{-1},\eeq which is a group homomorphism. %The same formula 
\source{very-stable-higgs-bundles-equivariant-multiplicity-mirror-symmetry}
\begin{lemma}
\source{very-stable-higgs-bundles-equivariant-multiplicity-mirror-symmetry}
\notready \label{pullnorm} For $M\in J$ we have $\pi^*(M)\in \tilde{J}$ and $$\tilde{\bP}_{\pi^*(M)}\cong {\mathrm{Nm}}_{\pi}^*(\bP_M).$$ 
\source{very-stable-higgs-bundles-equivariant-multiplicity-mirror-symmetry}
\end{lemma}
\begin{proof}  Fix $\tilde{c}_0\in \mathbf{C}$ a lift such that $c_0=\pi(\tilde{c}_0)$. Then by \cite[p. 331]{birkenhake-lange} the diagram $$\begin{array}{lcl} \Pic_0(\tilde{J}) &\stackrel{\alpha_{\tilde{c}_0}^*}{\rightarrow} & \tilde{J}\\  \!\!\!\!\!\!{{\mathrm{Nm}}}_{\pi}^*\uparrow && \uparrow  \pi^*\\ \Pic_0(J) &\stackrel{\alpha^*_{c_0}}{\rightarrow} & {J} \end{array}$$ commutes. Thus $$\alpha_{\tilde{c}_0}^*({\mathrm{Nm}}_{\pi}^*(\bP_M))=\pi^*(\alpha_c^*(\bP_M))=\pi^*(M).$$ By \eqref{char} this implies the statement. 
\end{proof}


\subsection{Fourier-Mukai transform of $\calO_{W_\calE^+}$}

In this section we will fix $d=-n(n-1)(g-1)$ the degree of our Higgs bundles. This will ensure that the line bundle $U$ defined in \eqref{definitionU} on a spectral curve corresponding to our Higgs bundle will have degree $0$. In particular, the canonical uniformising Higgs bundle $\calE_0\in \M$ and the canonical Hitchin section $W^+_{\calE_0}\subset \M$ will be contained in our Higgs moduli space $\M$.   

Consider now a very stable type $(1,\dots,1)$ Higgs bundle 
$\calE_\delta\in \M^{s\T}$. That means  $\delta=(\delta_0,\delta_1,\dots,\delta_{n-1}) \in J_\ell(C)\times C^{[m_1]}\times \dots \times C^{[m_{n-1}]}\subset \M$ such that $\calE_{\d}$ is very stable. %In other words $\delta_1+\dots + \delta_{n-1}$ is reduced, and we denote $\delta_i=c_{i1}+\dots+c_{im_i}$, for points $c_{ij}\in C$ or $-c_{0j}\in C$.





\subsubsection{Construction of $\Lambda_\delta$ from universal bundles}
\label{lambdadelta}
\source{very-stable-higgs-bundles-equivariant-multiplicity-mirror-symmetry}
The proposed mirror of $\calO_{W^+_{\calE_\delta}}$ will be a vector bundle $\Lambda_\delta$ on $\M^s$ constructed from universal Higgs bundles whose existence was proved in \cite{simpson}. We will have to overcome two difficulties in the construction. One is that the individual universal Higgs bundles will only be sheaves in a twisted sense, and they will become untwisted vector bundles only if an appropriate number of them is tensored together. Thus we will first construct universal bundles in an \'etale covering of $\M^s$ and only when an appropriate number will be tensored together the obstruction to glue will vanish. The second difficulty will be in that the universal Higgs bundles are only unique up to tensoring with line bundles pulled back from $\M^s$. Thus in order to construct well defined universal bundles we will have to normalise them appropriately. 


\'Etale locally $U_i\to \M^s$ we can define  \cite[Theorem 4.7.(4)]{simpson1} universal Higgs bundles $(\bE_{U_i},\bPhi_{U_i})$ on $U_i\times C$.  More precisely,  choosing a point $c_0\in C$, we use the framed moduli space $\calR^s$ \cite[Theorem 4.10]{simpson1}. Then  $C\times \calR^s$ carries  a unique universal framed Higgs bundle $(\bE_{\calR^s},\bPhi_{\calR^s},\bbbeta_{\calR^s})$, where the map $\bE_{\calR^s}$ is a rank $n$ vector bundle on $C\times \calR^s$, the universal Higgs field $\bPhi_{\calR^s}\in H^0(C\times\calR^s; \End(\bE_{\calR^s}\otimes K))$ and the universal framing $\bbbeta_{\calR^s}:\bE_{\calR^s}|_{\{c_0\}\times \calR^s}\stackrel{\cong}{\to} \calO_{\calR^s}^n$. We can take slices $U_i\subset \calR^s$ of the $\PGL_n$-principal bundle $\calR^s\to \M^s$ by  Luna's slice theorem \cite[Theorem III.1]{luna} such that $U_i\to \M^s$ is \'etale and $\coprod U_i \to \M^s$ covers $\M^s$. Then $(\bE_{U_i},\bPhi_{U_i})$ is just the Higgs bundle part of the family of framed Higgs bundles $(\bE_{U_i},\bPhi_{U_i},\bbbeta_{U_i})$ which is obtained as the pull back of the universal framed Higgs bundle $(\bE_{\calR^s},\bPhi_{\calR^s},\bbbeta_{\calR^s})$ from $C\times \calR^s$.  

The local universal framings $\bbbeta_{U_i}$  can be  used to normalise our local universal Higgs bundles. For this denote the {\em normalized determinant} $$\Nm(E):=\det(E)\otimes K^{n \choose 2}$$ for a vector bundle (or a family of vector bundles) on $C$. Fixing an isomorphism $K_{c_0}\cong \C$, a framing $\beta:E_{c_0}\to 
\C^n$ gives rise to a framing $\Nm(\beta):\Nm(E)|_{c_0}\cong \C$ of the line bundle $\Nm(E)$.  Then $(\Nm(\E_{U_i}),\Nm(\bbbeta_{U_i}))$ is a $U_i$-family of framed degree $0$ line bundles on $C$, thus unique from Remark~\ref{framedjacobian}. This gives our local normalisation  \beq \label{determinant}\phi_i:\Nm(\bE_{U_i})\stackrel{\cong}{\to}\Nm^*(\L),\eeq where $\L$ is the line bundle of \eqref{univlambda} and $\Nm:\M^s\to J$ is  given by \beq \label{normdet} \Nm(E,\Phi):=\Nm(E).\eeq  
\source{very-stable-higgs-bundles-equivariant-multiplicity-mirror-symmetry}






We will now attempt to glue our normalised local Higgs bundles together. On the overlap $U_i\times_{\M^s} U_j$  of two of our \'etale covers \cite[(4.2)]{hausel-thaddeus} implies that we get  a  line bundle $L_{ij}\in {\rm Pic}(U_i\times_{\M^s} U_j)$  which satisfies  $$(\bE_{U_i},\bPhi_{U_i})\cong (L_{ij}\bE_{U_j},\bPhi_{U_j}).$$ Passing to a refinement of the \'etale covering if necessary we can assume that all $L_{ij}$ are trivial. It is obvious how to do this in the analytic topology, the \'etale version which we need is  more complicated.   Namely, we can trivialise $L_{ij}$ on $U_i\times_{\M^s} U_j$ even on a Zariski covering, and then we can use \cite[Theorem 4.1]{artin} ( see also \cite[Lemma III.2.19]{milne}) to refine the original coverings so that the connecting line bundles will trivialise on the overlaps. 

This way we can assume that $$(\bE_{U_i},\bPhi_{U_i})\cong (\bE_{U_j},\bPhi_{U_j})$$ on $U_i\times_{\M^s} U_j$ for all ${ij}$. Taking such an isomorphism $$\phi_{ij}: \bE_{U_i} \stackrel{\cong}{\to} \bE_{U_j}$$ and recalling the fixed isomorphisms \eqref{determinant} we consider $$\phi_j\circ \Nm(\phi_{ij}) \circ \phi_i^{-1}:\Nm^*(\L)\to \Nm^*(\L).$$ By passing to an \'etale cover if needed we can assume that this function has an $n$th root (for this we use again \cite[Theorem 4.1]{artin}), and thus we can rescale $\phi_{ij}$ so that we can assume that $$\phi_j\circ \Nm(\phi_{ij}) \circ \phi_i^{-1}={\mathrm{Id}}_{\Nm^*(\L)}.$$ The isomorphisms $\phi_{ij}$ with this property are unique only up to multiplying with an $n$th root of unity, but we can make sure that $\phi_{ji}=\phi_{ij}^{-1}$.  
Thus $\phi_{jk}\circ\phi_{ij}$ agrees with $\phi_{ik}$ up to an $n$th root of unity, which we denote by $\theta_{ijk}\in \bbmu_n\subset \C^\times$ so that $$\phi_{jk}\circ\phi_{ij}=\theta_{ijk}\phi_{ik}.$$
We compute $$\theta_{ijk} \theta_{kli}{\mathrm{Id}}_{\bE_{U_k}}= \theta_{ijk}\phi_{ik}\circ \theta_{kli}\phi_{ki}=\phi_{jk}\circ\phi_{ij}\circ \phi_{li}\circ\phi_{kl}\in H^0(U_i\times_{\M^s}U_j\times_{\M^s}U_k\times_{\M^s}U_l;\End(\bE_{U_k}))$$ and similarly  
$$\theta_{jkl} \theta_{lij}Id_{\bE_{U_l}}= \phi_{kl}\circ\phi_{jk}\circ \phi_{ij}\circ\phi_{li}\in H^0(U_i\times_{\M^s}U_j\times_{\M^s}U_k\times_{\M^s}U_l;\End(\bE_{U_l})).$$ 
These two scalar endomorphisms are conjugate by $\phi_{kl}$ and therefore the scalars agree: $$\theta_{ijk} \theta_{kli}=\theta_{jkl} \theta_{lij}.$$ It follows that the $n$th roots of unity $\theta_{ijk}$ form a {\v{C}}ech $2$-cocycle $\theta$, thus define a cohomology class $[\theta]\in H^2(\M^s;\muhat_n)$, that is a $\muhat_n$-gerbe. By definition \cite[\S 1.2]{caldararu} the data of the vector bundles $\bE_{U_i}$ and  $\phi_{ij}$ constitute a $\theta$-twisted vector bundle $\bE$ on $\M^s\times C$, when $\theta$ is considered a  \v{C}ech $2$-cocycle for the sheaf $\calO_{\M^s}^*$ via the embedding $\bbmu_n\subset \C^\times$. We can multiply $\phi_{ij}$ by $\alpha_{ij}\in \muhat_n$ for a $1$-cocycle representing $[\alpha] \in H^1(\M^s; \bbmu_n)$, i.e. an order $n$ line bundle, then all of the above will hold. Thus the twisted bundle $\bE$ constructed above  is well-defined up to tensoring with an order $n$ line bundle on $\M^s$  but normalized such that \beq\label{normalise}\Nm(\bE)\cong \Nm^*(\L).\eeq We also note that the $2$-cocycle $\theta$ does not change if we rescale with the $1$- cocycle $\alpha$ as above. %Finally, for $N\in J$ a degree $0$ line bundle on $C$ from \eqref{translateL} we have \beq \label{translateE}\bE^{N^nM}=\nm^*(\bP_N)\otimes \bE.\eeq 
\source{very-stable-higgs-bundles-equivariant-multiplicity-mirror-symmetry}




For a rank $n$ vector bundle $E$ and line bundle $M$ we have  $\Lambda^{i}(ME)=M^i\Lambda^i(E)$. It follows that $\Lambda^i(\bE)$ is naturally a $\theta^i$-twisted sheaf. For  $(\ell,m_1,\dots,m_{n-1})\in \Z\times \Z_{\geq 0}$  and  $\d=(\delta_0,\delta_1,\dots,\delta_{n-1})\in J_\ell(C)\times C^{[m_1]}\times\dots\times C^{[m_{n-1}]}$  with $\delta_0=c_{01}+\dots+c_{0m_0}$ where $ c_{0j}\in C$ or $-c_{0j}\in C$ is a point in $C$ and  $\delta_i=c_{i1}+\dots + c_{im_i}\in C^{[m_i]} $ for points $c_{ij}\in C$  we consider~\footnote{It is easy to check \eqref{normdeterminant} that $\otimes_{j=1}^{m_0}\Lambda^n(\E_{c_{0j}})\cong \Nm^*\left(\bP_{\calO(\delta_0-\ell c_0)}\right)$ depends only on the line bundle $\calO(\delta_0)$ and not on the choice of divisor $\delta_0=c_{01}+\dots+c_{0m_0}$.}   \beq\label{bigprod}\Lambda_\delta:=\bigotimes_{i=0}^{n-1}\bigotimes_{j=1}^{m_i} \Lambda^{n-i}(\bE_{c_{ij}}), \eeq  where $\bE_{c_{ij}}=\bE|_{\M^s\times{\{c_{ij}\}}}$ when $c_{ij}\in C$ and $\bE_{c_{0j}}=\bE^*|_{\M^s\times{\{-c_{0j}\}}}$, when $-c_{0j}\in C$. Then for a line bundle $M$ on $\M^s$ we have \beq \label{twistedM}\bigotimes_{i=0}^{n-1}\bigotimes_{j=1}^{m_i} \Lambda^{n-i}(M\bE_{c_{ij}})=M^{D}\bigotimes_{i=0}^{n-1}\bigotimes_{j=1}^{m_i}  \Lambda^{n-i}(\bE_{c_{ij}}),\eeq where \beq \label{D}D=n\ell+\sum_{i=1}^n(n-i)m_i.\eeq Thus $\Lambda_\delta$ in \eqref{bigprod} is a $\theta^{D}$-twisted sheaf on $\M^s$. In particular, when $n|D$ the $2$-cocycle  $\theta^D=1$ is trivial, thus in this case \eqref{bigprod} descends as a vector bundle to $C\times\M^s$. Additionally, when $n|D$ the ambiguity of tensoring $\E$ with an order $n$ line bundle disappears, and so \eqref{bigprod} is a well-defined vector bundle on $\M^s$, which is normalised such that  \beq\label{bignorm}\det(\Lambda_\delta)=\det\left(\bigotimes_{i=0}^{n-1}\bigotimes_{j=1}^{m_i} \Lambda^{n-i}(\bE_{c_{ij}}))\right) \cong \Nm^*\left(\bigotimes_{i=0}^{n-1}\bigotimes_{j=1}^{m_i} \L_{c_{ij}}^{\bino{n-1}{n-i-1}}\right).\eeq 
\source{very-stable-higgs-bundles-equivariant-multiplicity-mirror-symmetry}

To study the dependence of the vector bundle $\Lambda_\delta$ on the base point $c_0\in C$ let  $c^\prime_0\in C$ be another point. The twisted vector bundle $\bE^\prime:=\bE \otimes \Nm^*(\pi_2^*\bP_{L})$, where $L^n\cong \calO(c_0-c_0^\prime)$, will be a twisted universal bundle normalised at $c^\prime_0$ in the sense that $\Nm(\bE^\prime)=\Nm^*(\L^\prime),$ where $\L^\prime$ is the universal line bundle \eqref{changenorm} normalised at $c^\prime_0$. When $D=0$, which we are assuming in this section, we see from \eqref{twistedM} that $\Lambda_\delta$ does not depend on the choice of the base point $c_0\in C$. 
 %Finally, for $N\in J$  $$\bigotimes_{i=1}^{n-1}\bigotimes_{j=1}^{m_i} \Lambda^{n-i}(\bE^{N^nM}_{c_{ij}})=\bigotimes_{i=1}^{n-1}\bigotimes_{j=1}^{m_i}  \Lambda^{n-i}(\nm^*{(\bP_N)}\otimes\bE_{c_{ij}})=\nm^*(\bP_{N^d})\otimes\bigotimes_{i=1}^{n-1}\bigotimes_{j=1}^{m_i}  \Lambda^{n-i}(\bE_{c_{ij}}),$$ thus \eqref{bigprod} does not change if we multiply $M$ by an order $d/n$ line bundle.  

\subsubsection{Fourier-Mukai transform over the good locus}

We shall compute the relative Fourier-Mukai transform of $\calO_{W_{\calE_\delta}^+}$ over a certain open subset of $\calA$. Denote by $\mathbf{C}\subset T^*C\times \calA$ the universal spectral curve. 
We denote by $\calA^{\#}\subset \calA$ the open subset of characteristic equations $a\in\calA^{\#}$ for which the following  properties hold \begin{enumerate} \item $C_a$ is smooth \item $\div(b)$ avoids the ramification divisor of the spectral cover $\pi_a:C_a\to C$  % \item  $W_\calE^+\cap h^{-1}(a)$ is transversal
	.\end{enumerate} We let $$\mathbf{C}^{\#}\subset T^*C\times \calA^{\#}$$ denote the family of spectral curves over $\calA^{\#}$. By the BNR correspondence we can identify $\M^{\#}:=h^{-1}(\calA^{\#})$ with the relative Jacobian of the family $\mathbf{C}^{\#}\to \calA^{\#}$: $$\M^{\#}\cong J({\mathbf{C}}^{\#}/{\calA}^{\#}).$$

We shall denote the relative Poincar\'e bundle on
$$J({\mathbf{C}}^{\#}/{\calA}^{\#})\times_{\calA^{\#}} J({\mathbf{C}}^{\#}/{\calA}^{\#})$$
by $\tilde{\bP}$. 
In this section
we will determine the relative Fourier-Mukai transform $$S(\calO_{W_{\calE_\delta}^+}|_{\M^\#})=(\pi_2)_*(\pi_1^*(\calO_{W_{\calE_\delta}^+}|_{\M^\#})\otimes^L \tilde{\bP})$$ over $\calA^{\#}$.

Recall that $\delta=(\delta_0,\delta_1,\dots,\delta_{n-1}) \in J_\ell(C)\times C^{[m_1]}\times \dots \times C^{[m_{n-1}]}\subset \M$ such that $\calE_{\d}$ is very stable. In other words $\delta_1+\dots + \delta_{n-1}$ is reduced, and we denote $\delta_i=c_{i1}+\dots+c_{im_i}$, for points $c_{ij}\in C$ or $-c_{0j}\in C$.  We also introduce the notation $$\calL_{\d}:=\calO_{W^+_{\calE_{\d}}}$$ for the structure sheaf of the Lagrangian upward flow of $\calE_{\d}$. We recall that in this section we set $d=-n(n-1)(g-1)$ for our Higgs bundles and clearly $d=D-n(n-1)(g-1)$ of \eqref{D} thus $D=0$ and 
$$\Lambda_{\d}:= \bigotimes_{i=0}^{n-1}\bigotimes_{j=1}^{m_i} \Lambda^{n-i}(\bE_{c_{ij}})$$ constructed in Section~\ref{lambdadelta} together with the normalisation \eqref{bignorm} is  a well-defined vector bundle on $\M^s$. With this notation we have

\begin{theorem}
\source{very-stable-higgs-bundles-equivariant-multiplicity-mirror-symmetry}
\notready \label{fouriermukai}   The relative Fourier-Mukai transform over 
\source{very-stable-higgs-bundles-equivariant-multiplicity-mirror-symmetry}
	$\calA^\#$  satisfies 
	\beq \label{mirrorl} S(\calL_{\d}|_{\M^\#})=\Lambda_{\d}|_{\M^\#}.\eeq
\source{very-stable-higgs-bundles-equivariant-multiplicity-mirror-symmetry}
\end{theorem}

%\begin{remark} As a sanity check we can see what this gives for $\delta=0$. Then $\calL_0=\calO_{W^+_0}$ is the structure sheaf of the Hitchin section. We expect its mirror to be the structure sheaf of the whole space $\M^s$. Indeed, choosing representative $0=c+-c$ for the trivial divisor on $C$ we get $\Lambda_0=\Lambda^n(\bE_{c})\otimes \Lambda^n(\bE_{-c})=\Lambda^n(\bE_{c})\otimes \Lambda^n(\bE^*_{c})\cong \calO_{\M^s}$. 
%\end{remark}

\begin{proof} First we  check the theorem fibrewise along $h$.    We let $a\in \calA^\#$ so that   $C_a$ is smooth, $\div(b)$ avoids the ramification divisor of the spectral cover $\pi_a:C_a\to C$ and  where $W_\calE^+\cap h^{-1}(a)$ is transversal.
	
	As $\calE_\delta$ is a very stable  type $(1,\dots,1)$ Higgs bundle, we have from Theorem~\ref{mainverystable}  that $b=b_{n-1}\circ \dots \circ b_1$ has no repeated $0$. We will denote by $c_{i1},\dots,c_{im_i}$ the $m_i$ zeroes of $b_i$. For $1\leq i\leq n-1$ and $1\leq j \leq m_i$  denote by $\tilde{c}^{1}_{ij},\dots,\tilde{c}^{n}_{ij}\subset C_a$ the $n$ points in the preimage $\pi_a^{-1}(c_{ij})$ of the zero $c_{ij}$ of $b_i$.
	From Proposition~\ref{vsintersect} we see that  $L\in W_\calE^+\cap h^{-1}(a)\cong W_\calE^+ \cap J(C_a)$ is a degree $0$ line bundle on $C_a$ of the form $L\cong \pi^*(E_1)(\Delta_1)$ where $\Delta_1=\sum_{i=1}^{n-1} \Delta_i-\Delta_{i+1}$. We know all possibilities for the divisor $\Delta_i-\Delta_{i+1}$ on $C_a$; namely it is effective and reduced and it has the form $$\Delta_i-\Delta_{i+1}=\sum_{j=1}^{m_i} \tilde{c}^{k^1_{ij}}_{ij}+\dots+\tilde{c}^{k^{n-i}_{ij}}_{ij}$$ where  $1\leq k^1_{ij}< \dots < k^{n-i}_{ij}\leq n$. Thus we see that \beq \label{sumsky}\calO_{W_\calE^+\cap h^{-1}(a)}=\bigoplus_{\{1\leq k^1_{ij}< \dots < k^{n-i}_{ij}\leq n\}^{1\leq j \leq m_i}_{1\leq i\leq n}} \calO_{\pi^*(E_1)\left(\sum_{i=1}^{n-1}\sum_{j=1}^{m_i} \tilde{c}^{k^1_{ij}}_{ij}+\dots+\tilde{c}^{k^{n-i}_{ij}}_{ij}\right)}\eeq is a direct sum of skyscraper sheaves. 
\source{very-stable-higgs-bundles-equivariant-multiplicity-mirror-symmetry}
	%Denoting $$D=\sum_i (n-i)m_i=-\deg(\pi^*(E_1))=-n\deg(E_1),$$ and 
	Fixing a lift $\tilde{c}_0\in C_a$ of $c_0\in C$ we can compute the Fourier-Mukai transform of \eqref{sumsky} as follows:
	\bes S(\calO_{W_\calE^+\cap h^{-1}(a)})&=& \bigoplus_{\{1\leq k^1_{ij}< \dots < k^{n-i}_{ij}\leq n\}^{1\leq j \leq m_i}_{1\leq i\leq n}} S\left(\calO_{\pi^*(E_1)\left(\sum_{i=1}^{n-1}\sum_{j=1}^{m_i} \tilde{c}^{k^1_{ij}}_{ij}+\dots+\tilde{c}^{k^{n-i}_{ij}}_{ij}\right)}\right)\\
	&=& \bigoplus_{\{1\leq k^1_{ij}< \dots < k^{n-i}_{ij}\leq n\}^{1\leq j \leq m_i}_{1\leq i\leq n}} S\left(\calO_{\pi^*(E_1)\left((\sum_{i=1}^n(n-i)m_i)\tilde{c}_0+\sum_{i=1}^{n-1}\sum_{j=1}^{m_i} \tilde{c}^{k^1_{ij}}_{ij}-\tilde{c}_0+\dots+\tilde{c}^{k^{n-i}_{ij}}_{ij}-\tilde{c}_0\right)}\right)\\ &=&\bigoplus_{\{1\leq k^1_{ij}< \dots < k^{n-i}_{ij}\leq n\}^{1\leq j \leq m_i}_{1\leq i\leq n}} \tilde{\P}_{\pi^*(E_1)(-n\ell \tilde{c}_0)}\otimes  \bigotimes_{i=1}^{n-1}\bigotimes_{j=1}^{m_i} \tilde{\P}_{( \tilde{c}^{k^1_{ij}}_{ij}-\tilde{c}_0)}\otimes \dots \otimes \tilde{\P}_{(\tilde{c}^{k^{n-i}_{ij}}_{ij}-\tilde{c}_0)}\\
	&=&  \tilde{\P}_{\pi^*(E_1)(-n\ell \tilde{c}_0)}\otimes \bigotimes_{i=1}^{n-1}\bigotimes_{j=1}^{m_i} \Lambda^{n-i} \left(\tilde{\P}_{(\tilde{c}^{1}_{ij}-\tilde{c}_0)}\oplus \dots \oplus \tilde{\P}_{(\tilde{c}^{n}_{ij}-\tilde{c}_0)} \right) 
	\ees
	
	We see that $$\tilde{\P}_{(\tilde{c}^{k}_{ij}-\tilde{c}_0)}=\tilde{\L}|_{J(C_a)\times \{\tilde{c}^{k}_{ij}\}},$$ where $\tilde{\L}$ is a universal Poincar\'e bundle on $J(C_a)\times C_a$ normalized at $\tilde{c}_0$ as defined in \eqref{univlambda}. Furthermore, \beq\label{universal}\tilde{\P}_{(\tilde{c}^{1}_{ij}-\tilde{c}_0)}\oplus \dots \oplus \tilde{\P}_{(\tilde{c}^{n}_{ij}-\tilde{c}_0)}=\pi_*(\tilde{\L})|_{h^{-1}(a)\times\{c_{ij}\}},\eeq and $(\pi_*(\tilde{\L}),\bPhi_{h^{-1}(a)})$ with $\bPhi_{h^{-1}(a)}:=\pi_*(x:\tilde{\L}\to \tilde{\L}\otimes \pi^*(K))$ is a universal Higgs bundle on $h^{-1}(a)\times C$. 
\source{very-stable-higgs-bundles-equivariant-multiplicity-mirror-symmetry}
	
	\begin{lemma}
\source{very-stable-higgs-bundles-equivariant-multiplicity-mirror-symmetry}
\notready\label{detpull} For  $L\in J(C_a)$ let ${\mathrm{Nm}}:J(C_a)\to J(C)$ given by ${\mathrm{Nm}}(L)=\det(\pi_*(L))\otimes K^{n\choose 2}$ the usual norm map. Then on $J(C_a)\times C$  we have $$\Nm((\pi_a)_*(\tilde{\L}))={\mathrm{Nm}}^*(\L),$$ where $\tilde{\L}$ is a universal line bundle on $J(C_a)\times \tilde{C_a}$ normalized at $\tilde{c}_0\in \tilde{C_a}$ and $\L$ is a universal line bundle on $J(C)\times C$ normalized at $c_0=\pi_a(\tilde{c}_0)$. 
\source{very-stable-higgs-bundles-equivariant-multiplicity-mirror-symmetry}
	\end{lemma}
	\begin{proof} For  $L\in J(C_a)$ we have $${\mathrm{Nm}}^*(\L)|_{\{L\}\times C}={\mathrm{Nm}}(L)=\Nm(\pi_*(\tilde{\L}))|_{\{L\}\times C}
		.$$
		
		On the other hand for $c\in C$ we have $${\mathrm{Nm}}^*(\L)|_{J(C_a)\times\{c\} }={\mathrm{Nm}}^*(\L|_{J(C)\times\{c\}})={\mathrm{Nm}}^*(\P_{(c-c_0)}).$$
		While for the other side $$\Nm(\pi_*(\tilde{\L}))|_{J(C_a)\times\{c\} }=\det(\pi_*(\tilde{\L})|_{J(C_a)\times\{c\} })=\det\left(\tilde{\P}_{(\tilde{c}_1-\tilde{c}_0)}+\dots+\tilde{\P}_{(\tilde{c}_n-\tilde{c}_0)}\right)=\tilde{\P}_{\pi^*(c-c_0)}$$ where $\pi^{-1}(c)=\{\tilde{c}_1,\dots,\tilde{c}_n\}$.
		
		Now Lemma~\ref{pullnorm} implies that  $$\det(\pi_*(\tilde{\L}))|_{J(C_a)\times\{c\} }= {\rm det}^*(\L)|_{J(C_a)\times\{c\} }.$$
		
		The result follows. 
	\end{proof}
	The norm map ${\mathrm{Nm}}:J(C_a)\to J(C)$ in Lemma~\ref{detpull} agrees with the one restricted from
	the normalised determinant map in \eqref{normdet} to $h^{-1}(a)\cong J(C_a)$. Thus $\Nm(\pi_*(\tilde{\L}))=\Nm^*({\L})$ and therefore  $\pi_*(\tilde{\L})\cong \bE$ on $h^{-1}(a)$. Furthermore $$ \tilde{\P}_{\pi^*(E_1)(-n\ell \tilde{c}_0)}\cong \Nm^*(\P_{E_1(-\ell c_0)})$$ from Lemma~\ref{pullnorm}. By the normalisation \eqref{normalise}, the definition \eqref{univlambda} and $E_1\cong\calO(\delta_0)=\calO(c_{01}+\dots+c_{0m_0})$   we get \beq\label{normdeterminant}\Nm^*(\P_{E_1(-\ell c_0)})\cong \bigotimes_{j=1}^{m_0} \Lambda^{n}(\bE_{c_{0j}}).\eeq Thus \eqref{universal} and Lemma~\ref{detpull} verify Theorem~\ref{universal} on $h^{-1}(a)$.
\source{very-stable-higgs-bundles-equivariant-multiplicity-mirror-symmetry}
	
	By noting that   \'etale-locally on $\calA^\#$  we have sections of the universal spectral curve corresponding to the preimages of zeroes of $b$ and $c_0$, we see that the same argument can be repeated relative to \'etale neighbourhoods $U_i\to \calA^\#$, proving Theorem~\ref{fouriermukai} \'etale-locally, which as argued above glue together to prove Theorem~\ref{fouriermukai}.
\end{proof}

\begin{remark}
\source{very-stable-higgs-bundles-equivariant-multiplicity-mirror-symmetry}
\notready As a sanity check we can see what this gives for $\delta=0$. Then $\calL_0=\calO_{W^+_0}$ is the structure sheaf of the canonical Hitchin section. We expect its mirror to be the structure sheaf of the whole space $\M^s$. Indeed, choosing the representative $0=c+-c$ for the trivial divisor on $C$ we get $\Lambda_0=\Lambda^n(\bE_{c})\otimes \Lambda^n(\bE_{-c})=\Lambda^n(\bE_{c})\otimes \Lambda^n(\bE^*_{c})\cong \calO_{\M^s}$. 
	
	More generally, we can determine $\Lambda_\delta$ in the case when $\delta_1=\dots=\delta_{n-1}=0$ and $\ell=0$, i.e. when $\calE_\delta$ is a uniformising Higgs bundle and $\calL_\delta=W^+_\delta$ is the structure sheaf of a Hitchin section. We see from \eqref{normdeterminant} that $\Lambda_\delta=\Nm^*(\bP_{\calO(\delta_0)})$ is a line bundle on $\M^s$ pulled back from the Jacobian $J$. 
\end{remark}

\begin{remark}
\source{very-stable-higgs-bundles-equivariant-multiplicity-mirror-symmetry}
\notready \label{remarkdual} Vector bundle duality defines a natural  involution \beq \label{iota} \begin{array}{cccc} \iota:&\M&\to &\M\\ &(E,\Phi)&\mapsto &(E^*K^{1-n},\Phi^T)\end{array}.\eeq where $\Phi^T$ is the induced action  of $\Phi$ on $E^*$. This is well-defined as $\deg(E^*K^{1-n})=-\deg(E)-n(n-1)2(g-1)=-n(n-1)(g-1)$. From the exact sequence \eqref{exactu} we see that $(\pi_a)_*(U^*) =\iota(E,\Phi)$ and therefore $\iota$ restricts to the abelian variety $h^{-1}(a)\cong J(C_a)$ as the inverse map $-1$. If $\calE\in \M$ we will denote $\calE^\iota:=\iota(\calE)$, then $\iota(W^+_\calE)=W^+_{\calE^\iota}$.  
\source{very-stable-higgs-bundles-equivariant-multiplicity-mirror-symmetry}
	
	For example for a type $(1,\dots,1)$ very stable Higgs bundle $\calE_{\d}$ we can determine $\calE^\iota_{\d}=\calE_{ \d^\iota}$, where  $\d^\iota=(\delta_0^\iota,\delta_1^\iota,\dots,\delta_{n-1}^\iota)=(-\delta_0-\delta_1-\dots-\delta_{n-1},\delta_{n-1},\dots,\delta_1)$.  We thus see that the skew involutive property \eqref{skewinv} of the Fourier-Mukai transform  yields
	\beq \label{mirrormirror} S(\Lambda_{\d}|_{\M^\#})=\calL_{\d^\iota}[-n^2(g-1)-1]|_{\M^\#}.\eeq 
\source{very-stable-higgs-bundles-equivariant-multiplicity-mirror-symmetry}
\end{remark}
% \begin{remark} The other identity \eqref{skewdual} from \cite{mukai} shows that a self-dual coherent sheaf $\calF$, i.e. $\calF^\vee\cong \calF$ on an Abelian variety is mapped to a skew self-dual coherent sheaf i.e.  $(-1)^*S(\calF)^\vee=S(\calF)$.

%\end{remark}

\subsection{An agreement for $\T$-equivariant Euler forms} 

As a motivation recall that mirror symmetry should be an equivalence of categories. Not only objects should match but morphisms between them. If two objects have $\T$-equivariant structure so will the vector space of morphisms between them. Thus its $\T$-character should agree with the $\T$-character of the vector space of morphisms between the mirror objects. In this subsection we find the agreement of these $\T$-characters (aka Euler pairing) for our mirror pairs of branes. Incidently, it is easier to pair a Lagrangian brane with a hyperholomorphic one because the computation reduces to a computation on the Lagrangian support which in our case is just a vector space. 


\subsubsection{$\T$-equivariant Euler characteristics and Euler forms}
\label{equivarianteuler}
\source{very-stable-higgs-bundles-equivariant-multiplicity-mirror-symmetry}
Recall that for a $\T$-equivariant coherent sheaf $\calF$ on a semiprojective variety $M$ the {\em equivariant Euler characteristic} is defined by
$$\chi_\T(M;\calF):=\sum_{ij} (-1)^i \dim(H^i(M;\calF)^{j}) t^{-j},$$ 
when the weight spaces $H^i(M;\calF)^{j}$, where $\T$ acts with weight $j$, are finite dimensional. The finiteness condition holds for smooth semiprojective varieties and  then $\chi_\T(M;\calF)\in \Z((t)),$ i.e.  $H^i(M;\calF)^j=0$ for $i$ sufficiently large. More generally for a $\T$-equivariant bounded complex of coherent sheaves $\calF^\bullet$ we define $\chi_\T(M;\calF^\bullet)=(-1)^k\sum_k \chi_\T(M;\calF^k)$. 

Furthermore, given two $\T$-equivariant coherent sheaves $\calF_1$ and $\calF_2$ on $M$ we define the {\em equivariant Euler pairing} as

\begin{multline*}\chi_\T(M;\calF_1,\calF_2)=\sum_{k,l} \dim(H^k({ R} {\mathcal Hom}(\calF_1,\calF_2))^l) (-1)^k t^{-l}=\\ =\sum_{k,l} \dim\left(\Hom_{{\mathbf D}_{coh}(M)}(\calF_1,\calF_2[k])^l\right) (-1)^k t^{-l}=\sum_{k,l} \dim(\Ext^k(M;\calF_1,\calF_2)^l) (-1)^k t^{-l} .\end{multline*}

As $\chi_\T(M;\calF_1,\calF_2)=\chi_\T(M;\calF_1^\vee\otimes^L\calF_2),$ where $\calF_1^\vee$ is the derived dual and $\otimes^L$ is the derived tensor product, we expect that  $\chi_\T(M;\calF_1,\calF_2)\in \Z((t))$ on a semiprojective variety $M$.

We note that the upward flow $W^+_\calE$ for $\calE\in \M^{s\T}$ is invariant under the $\T$-action and we shall endow $\calO_{W^+_\calE}$  with the trivial $\T$-equivariant structure. 
\subsubsection{$\T$-equivariant universal bundles}
\label{equiuni}
\source{very-stable-higgs-bundles-equivariant-multiplicity-mirror-symmetry}
On the other hand we would like to endow the  sheaf \eqref{bigprod}, when $n|D$,  on $\M^s\times C$ with a $\T$-equivariant structure. For this we need

\begin{lemma}
\source{very-stable-higgs-bundles-equivariant-multiplicity-mirror-symmetry}
\notready \label{teqcover} We can choose the \'etale cover $\coprod_i U_i\to \M^s$ so that the $U_i$ are $\T$-equivariant and  the local universal Higgs bundle  $(E_{U_i},\bPhi_{U_i})$ on $U_i\times C$ carries a $\T$-equivariant structure with $\T$ acting on $\bPhi_{U_i}$ with weight $-1$. Furthermore the isomorphism \eqref{determinant} is $\T$-equivariant with some $\T$-equivariant structure on $\L$.
\source{very-stable-higgs-bundles-equivariant-multiplicity-mirror-symmetry}
\end{lemma}

\begin{proof}
	First we note that over $\calR\times C$ there exists, from \cite[Theorem 4.10]{simpson1}, the universal  framed Higgs bundle $(\bE_\calR,\bPhi_\calR,{\bbbeta}_\calR)$. We have the $\T$-action on $\calR$ multiplying the Higgs field and we  denote  the morphism $m_\lambda:\calR\to \calR$ given by scaling the Higgs field with $\lambda\in \T$. Then $$(m_\lambda^*(\bE_\calR),\lambda^{-1}m_\lambda^*(\bPhi_\calR),m_\lambda^*(\bbbeta_\calR))$$ is also a universal framed Higgs bundle and as $\calR$ is a fine moduli space with a unique universal object we get that $\bE_\calR\cong m_\lambda^*(\bE_\calR)$ canonically, inducing $f_\lambda:\bE_\calR\to \bE_\calR$ covering $m_\lambda$ and making the diagram \beq\label{univdiag}\begin{array}{ccc} \bE_\calR & \stackrel{\bPhi_\calR}{\longrightarrow} &\bE_\calR\otimes K \\ f_\lambda \downarrow && \downarrow f_\lambda  \\ \bE_\calR & \stackrel{\lambda^{-1} \bPhi_\calR}{\longrightarrow} &\bE_\calR\otimes K \end{array}\eeq commute. This will endow $\bE_\calR$ with a canonical $\T$-equivariant structure which acts on $\bPhi_\calR$ with weight $-1$.
\source{very-stable-higgs-bundles-equivariant-multiplicity-mirror-symmetry}
	
	Now we take a  point $\calE\in \M^{s}$ such that the orbit $\T\calE\subset \M^s$ is closed. This means that  either  $\calE\in \M^{s\T}$ or that  neither $\lim_{\lambda\to 0}\lambda\calE$ nor (if it exists) $\lim_{\lambda\to \infty}\lambda\calE$ is stable. We take  a lift $\tilde{\calE}\in \calR^s$. Recall that $\PGL_n$ acts freely on $\calR^s$ by changing the framing and $\T$ acts by multiplying the Higgs field. The $\PGL_n$ quotient $p:\calR^s\to \M^s$ is $\T$-equivariant. It follows  that for $\G:=\T\times \PGL_n$ the orbit $\G\tilde{\calE}\subset \calR^s$ is closed. We let $\G_{\tilde{\calE}}\subset \G$ denote the stabilizer which projects injectively into $\T$ as $\PGL_n$ acts freely.  We apply Luna's slice theorem \cite[Theorem III.1]{luna} for the action of $\G$ on $\calR^s$ around the point $\tilde{\calE}$. This way we get a locally closed $\G_{\tilde{\calE}}$-invariant affine slice $\tilde{\calE} \in V\subset \calR^s$, such that the map  $\G\times_{\G_{\tilde{\calE}}} V \to \calR^s$ is \'etale and whose image $U\subset \calR^s$ is  a saturated affine open $\G$-invariant subset, the morphism \beq \label{etale}V/\G_{\tilde{\calE}} \to U/\G\eeq is also \'etale and the natural morphism induces \beq \label{iso}\G\times_{\G_{\tilde{\calE}}} V \cong U\times_{U/\G} V/\G_{\tilde{\calE}}\eeq a $\G$-isomorphism.  We let $V^\prime:=\T\times_{\G_{\tilde{\calE}}}V=(\G\times_{\G_{\tilde{\calE}}}V)/\PGL_n$ then $V^\prime/\T=V/\G_{\tilde{\calE}}$ and dividing \eqref{iso} by $\PGL_n$ gives $$V^\prime=\T\times_{\G_{\tilde{\calE}}} V\cong U/\PGL_n\times_{U/\G} V/\G_{\tilde{\calE}}.$$ By base change of the \'etale morphism \eqref{etale}  by $U/{\PGL_n}\to U/\G=(U/\PGL_n)/\T$  \beq\label{etalelocal}V^\prime \to U/\PGL_n\subset \M^s\eeq is also \'etale and $\T$-equivariant. Finally, the $\T$-equivariant universal Higgs bundle $(\bE_\calR,\bPhi_{\calR})$ restricts to the $\T$-equivariant $V^\prime\to  \calR$ giving  $\T$-equivariant universal bundles $(\bE_{V^\prime},\bPhi_{V^\prime})$ \'etale locally \eqref{etalelocal} around $\calE\in U/\PGL_n$. 
\source{very-stable-higgs-bundles-equivariant-multiplicity-mirror-symmetry}
	
	Finally, we note that a $\T$-invariant open covering of the closed $\T$-orbits in $\M^s$ covers the whole $\M^s$ and so the result follows.  
	
\end{proof}

Thus we have local $\T$-equivariant universal Higgs bundles on $U_k\to \M^s$. On overlaps $C\times U_k\times_{\M^s} U_l$ they differ by a $\T$-equivariant line bundle by \cite[(4.2)]{hausel-thaddeus}. Thus the $\T$-equivariant projective bundles $\bP\E_{U_i}$ and $\bP\E_{U_j}$ are canonically isomorphic over the overlap. Therefore they descend as a $\T$-equivariant projective bundle $\bP\E$ over $C\times \M^s$. 

\begin{lemma}
\source{very-stable-higgs-bundles-equivariant-multiplicity-mirror-symmetry}
\notready In the notation of \eqref{bigprod} when $n|D=n\ell+\sum_{i=1}^{n-1}(n-i)m_i$ there is a $\T$-equivariant structure on the vector bundle $\Lambda_\d=\bigotimes_{i=0}^{n-1}\bigotimes_{j=1}^{m_i} \Lambda^{n-i}(\bE_{c_{ij}})$ normalised such that 
	\beq\label{equinormalise} \det\left(\bigotimes_{i=0}^{n-1}\bigotimes_{j=1}^{m_i} \Lambda^{n-i}(\bE_{c_{ij}}))\right)\cong \Nm^*\left(\bigotimes_{i=0}^{n-1}\bigotimes_{j=1}^{m_i} \L_{c_{ij}}^{\bino{n-1}{n-i-1}}\right)\eeq as equivariant line bundles, where the line bundle $\L_{c_{ij}}^{\bino{n-1}{n-i-1}}$ on $J$ is endowed with a weight ${n-i\choose 2}-{n\choose 2}{\bino{n-1}{n-i-1}}$ $\T$-action.
\source{very-stable-higgs-bundles-equivariant-multiplicity-mirror-symmetry}
\end{lemma}

\begin{proof}
	Similarly to $\P\bE$ above,  we can argue that $\bigotimes_{i=0}^{n-1}\bigotimes_{j=1}^{m_i} \Lambda^{n-i}(\bE_{U_k}|_{c_{ij}\times U_k})$ on $U_k\to \M^s$  carries a $\T$-equivariant structure, and they differ by a $\T$- equivariant line bundle on the overlaps. Thus the total space of the projective bundle $\P\left(\bigotimes_{i=0}^{n-1}\bigotimes_{j=1}^{m_i} \Lambda^{n-i}(\bE_{c_{ij}})\right)$ descends to $\M^s$ with a $\T$-equivariant structure.  When $n|D=\sum_{i=0}^{n-1}(n-i)m_i$ then we saw before Theorem~\ref{fouriermukai} that even the vector bundle $\bigotimes_{i=0}^{n-1}\bigotimes_{j=1}^{m_i} \Lambda^{n-i}(\bE_{c_{ij}})$ descends to $\M^s$. Thus the projective bundle $\P\left(\bigotimes_{i=0}^{n-1}\bigotimes_{j=1}^{m_i} \Lambda^{n-i}(\bE_{c_{ij}})\right)$ has a sheaf $\calO(1)$ on it, such that \beq \label{pushdown} \pi_*(\calO(1))^*\cong \bigotimes_{i=0}^{n-1}\bigotimes_{j=1}^{m_i} \Lambda^{n-i}(\bE_{c_{ij}}).\eeq By \cite[Theorem 4.2.2]{brion} the $\T$-action on  $\P\left(\bigotimes_{i=0}^{n-1}\bigotimes_{j=1}^{m_i} \Lambda^{n-i}(\bE_{c_{ij}})\right)$ lifts to $\calO(1)$ because $\Pic(\T)$ is trivial and thus the last map in \cite[(12)]{brion} is an isomorphism. Thus we have a $\T$-equivariant structure on $\bigotimes_{i=0}^{n-1}\bigotimes_{j=1}^{m_i} \Lambda^{n-i}(\bE_{c_{ij}})$ from \eqref{pushdown}. 
\source{very-stable-higgs-bundles-equivariant-multiplicity-mirror-symmetry}
	
	We can modify this $\T$-equivariant structure by tensoring with a $\T$ -equivariant  structure on $\calO_{\M^s}$. We can choose it  such that at the canonical uniformising  Higgs bundle $\calE_0=\calE_{\calO_C,0}$  we get that
	$\bigotimes_{i=0}^{n-1}\bigotimes_{j=1}^{m_i} \Lambda^{n-i}(\bE_{c_{ij}}))|_{\calE_0}$ has only non-positive weights and non-trivial weight $0$ space. We can endow $E_{\calO_C}=\calO\oplus K^{-1}\oplus \dots \oplus K^{-n+1}$ with a $\T$-action which acts by weight $-i$ on the $i$th summand $K^{-i}$. This will act on $\Phi_{0}$ with weight $-1$. Therefore any $\T$-equivariant family of stable Higgs bundles containing $\calE_0$ will restrict to $\calE_0$ as this $\T$-equivariant structure, up to a rescaling of the $\T$-action by an overall $1$-dimensional $\T$-module. Therefore  our normalisation of the $\T$-action on  $\bigotimes_{i=0}^{n-1}\bigotimes_{j=1}^{m_i} \Lambda^{n-i}(\bE_{c_{ij}})$  will satisfy \beq\label{eqmultuniv}\chi_\T(\bigotimes_{i=0}^{n-1}\bigotimes_{j=1}^{m_i} \Lambda^{n-i}(\bE_{c_{ij}}))|_{\calE_0})=\prod_{i=0}^{n-1} \left[\begin{array}{c} n \\ n-i\end{array}\right]^{m_i}_t\eeq which incidently also agrees with the equivariant multiplicity $m_{\calE_{\d}}(t)$ of \eqref{explicit}. This will be related to  mirror symmetry in \eqref{mirroriden}. 
\source{very-stable-higgs-bundles-equivariant-multiplicity-mirror-symmetry}
	
	A straightforward computation checks that this normalisation of the $\T$-action will induce the one on the determinant as claimed in the Lemma. 
\end{proof}
%For one of our $\T$-equivariant \'etale covers $U_k\to \M^s$ from Lemma~\ref{teqcover} we will get a $\T$-equivariant structure on $\bigotimes_{i=1}^{n-1}\bigotimes_{j=1}^{m_i} \Lambda^{n-i}(\bE_{U_k}|_{c_{ij}\times U_k})$

% A simple computation gives  $$\det(\bigotimes_{i=1}^{n-1}\bigotimes_{j=1}^{m_i} \Lambda^{n-i}(\bE_{U_k}))= \det(\E_{U_i})^{rD/n},$$ where $r=\rank(\det(\bigotimes_{i=1}^{n-1}\bigotimes_{j=1}^{m_i} \Lambda^{n-i}(\bE_{U_k})))$ and $D$ is qiven by \eqref{D}.
%Note that $\nm:\M^s\to J$ is $\T$-equivariant with the trivial $\T$-action on $J$, thus the $\T$-equivariant line bundle $\nm^*(\L)$ 
%from the trivial $\T$-equivariant structure on $\L$ on $C\times J$, will differ from $\det(\E_{U_k})$ by a trivial line bundle with a $\T$-equivariant structure on it. We have \cite[Lemma 4.2.1]{brion} that such a $\T$-equivariant  line bundle is $L_{\chi_k}$ induced by a character $\chi_k:\C^\times \to \C^\times$. Thus as $\T$-equivariant line bundles $\det(\E_{U_k})=L_{\chi_k} \nm^*(\L)$.  Consequently $\det(\E_{U_k|_{c_{ij}\times U_k}})=L_{\chi_k} \nm^*(\L)|_{c_{ij}\times U_k}$. We can thus compute 
% $$\det(\bigotimes_{i=1}^{n-1}\bigotimes_{j=1}^{m_i} \Lambda^{n-i}(\bE_{U_k}|_{c_{ij}\times U_k}))=\bigotimes_{i=1}^{n-1}\bigotimes_{j=1}^{m_i} \det(\bE_{U_k}|_{c_{ij}\times U_k})^{r(n-i)/n}=L_{\chi_k}^{rD/n}\bigotimes_{i=1}^{n-1}\bigotimes_{j=1}^{m_i} \nm^*(\L)|_{c_{ij}\times U_k}, $$
% where $r=\rank(\bigotimes_{i=1}^{n-1}\bigotimes_{j=1}^{m_i} \Lambda^{n-i}(\bE_{U_k}|_{c_{ij}\times U_k}))$. Thus when $m=D/n$ an integer we can modify the $\T$ structure on our bundle so that $$\det(L_{\chi_k}^{-m}\bigotimes_{i=1}^{n-1}\bigotimes_{j=1}^{m_i} \Lambda^{n-i}(\bE_{U_k}|_{c_{ij}\times U_k}))=\bigotimes_{i=1}^{n-1}\bigotimes_{j=1}^{m_i} \nm^*(\L)|_{c_{ij}\times U_k}, $$

%\det(\L)_{c_{ij}\times U_k}



%We have checked in Theorem~\ref{fouriermukai} that fiberwise Fourier-Mukai transform on generic fibers of the Hitchin map takes $\calO_{W^+_\calE}$, the structure sheaf of the upward flow of very stable type $(1,\dots,1)$ Higgs bundle to  a certain  tensor product of exterior powers of the universal bundle at points in $C$. Based on the  mirror symmetry proposals \cite{hausel-thaddeus,kapustin-witten,donagi-pantev} we also expect that the mirror of $\calO_{W^+_\calE}$ the upward flow should agree with the same tensor product of universal bundles (at least over the stable locus). As mirror symmetry should induce an equivalence of derived categories of coherent sheaves, it should preserve the Euler pairing, and the following can be considered a numerical consequence of our mirror symmetry proposal. 

  The following is an indication  that the mirror relationships \eqref{mirror} and \eqref{mirrormirror}  over $\calA^\#$ hold over the whole $\calA$. 

\begin{theorem}
\source{very-stable-higgs-bundles-equivariant-multiplicity-mirror-symmetry}
\notready \label{euleriso} For $\ddelta=(\delta_0,\delta_1,\dots,\delta_{n-1})\in J_{{\ell}}(C)\times C^{[m_1]}\times \dots \times C^{[m_{n-1}]}$ and $\ddelta^\prime=({\delta_0}^\prime,\delta^\prime_1,\dots,\delta^\prime_{n-1})\in J_{{\ell^\prime}}(C)\times C^{[m^\prime_1]}\times \dots \times C^{[m^\prime_{n-1}]}$  let $\calE_{\ddelta}, \calE_{\ddelta^\prime}\in \M^{s\T}$ be two very stable type $(1,\dots,1)$ Higgs bundles. Then \bes \chi_\T\left(\M;\calL_{\ddelta^\prime},\Lambda_{\ddelta}\right)&=&\chi_\T(\det(\C_{1}\otimes\calA^*)) \chi_\T\left(\M;\Lambda_{\ddelta^\prime},\calL_{\ddelta^\iota}[-n^2(g-1)-1]
\source{very-stable-higgs-bundles-equivariant-multiplicity-mirror-symmetry}
	\right)
	\\ &= &(-1)^{n^2(g-1)+1} t^{(4n+1)(n-1)n(g-1)/6}  \chi_\T\left(\M;\Lambda_{\ddelta^\prime},\calL_{\ddelta^\iota}\right).\ees
	
\end{theorem}   	\begin{proof} We compute
	
	\begin{multline*} \chi_\T\left(\M;\calL_{\ddelta^\prime},\Lambda_{\ddelta}\right)= \chi_\T\left(\M;\calL_{\ddelta^\prime}^\vee \otimes ^L \Lambda_{\ddelta}\right)=\\ = \chi_\T\left(W^+_{\d^\prime};\det(N_{W^+_{\d^\prime}})[-n^2(g-1)-1] \otimes \Lambda_{\ddelta}|_{W^+_{\d^\prime}}\right),\end{multline*} where in the last equation we used \cite[Corollary 3.40]{huybrechts} and the simplified notation $W^+_{\d^\prime}:=W^+_{\calE_{\d^\prime}}$.
	
	As $\det(N_{W^+_{\d^\prime}}) \otimes \Lambda_{\ddelta}|_{W^+_{\d^\prime}}$ is a $\T$-equivariant vector bundle on the  affine space $W^+_{\d^\prime}$ its $\T$-equivariant index can be computed as $$\chi_\T\left(W^+_{\d^\prime};\det(N_{W^+_{\d^\prime}}) \otimes \Lambda_{\ddelta}|_{W^+_{\d^\prime}}\right)=\chi_\T\left((\det(N_{W^+_{\d^\prime}}) \otimes \Lambda_{\ddelta})|_{{\calE_{\d^\prime}}}\right) \chi_\T(W^+_{\d^\prime};\calO_{W^+_{\d^\prime}}).$$
	
	The same argument as above for \eqref{eqmultuniv} gives  $$\chi_\T(\Lambda_{\d}|_{\calE_{\d^\prime}})=m_{\calE_{\d}}(t).$$ From Definition~\ref{equmult} we have $$\chi_\T(W^+_{\d^\prime};\calO_{W^+_{\d^\prime}})=m_{\calE_{\d^\prime}}(t)\chi_\T(\Sym(\calA^*)).$$ Finally $N_{W^+_{\d^\prime}}|_{\calE_{\d^\prime}}\cong T^{\leq  0} _{\calE_{\d^\prime}}$ %and our homogeneity $1$ symplectic form $\omega$ gives an isomorphism $ T^{-*} _{\calE_{\d^\prime}}\cong \C_{-1}\otimes T^+_{\calE_{\d^\prime}}$ of $\T$-modules where $\C_{k}$ denotes the character of weight $k\in \Z$.  
	consequently $$\chi_\T(\det(N_{W^+_{\d^\prime}})|_{{\calE_{\d^\prime}}}))= \chi_\T(\det(T^{\leq  0}_{\calE_{\d^\prime}}))$$  as $\dim \calA = \dim T^+_{\calE_{\d^\prime}}=\dim \M/2$. Thus for the left hand side we have 
	\beq \label{lhs}	\chi_\T\left(\M;\calL_{\ddelta^\prime},\Lambda_{\ddelta}\right)=(-1)^{\dim \cal A}m_{\calE_{\d}}(t)m_{\calE_{\d^\prime}}(t)\chi_\T(\Sym(\calA^*))\chi_\T(\det(T^{\leq  0}_{\calE_{\d^\prime}})).\eeq
\source{very-stable-higgs-bundles-equivariant-multiplicity-mirror-symmetry}
	
	For the other side we compute similarly \beq \chi_\T\left(\M;\Lambda_{\ddelta^\prime},\calL_{\ddelta^\iota}\right)\nonumber&=&\chi_\T(\M;\Lambda^\vee_{\ddelta^\prime}\otimes^L\calL_{\ddelta^\iota})\\ \nonumber&=&\chi_\T(W^+_{\d^\iota};\Lambda^\vee_{\ddelta^\prime}|_{W^+_{\d^\iota}} )\\ \nonumber&=&\chi_\T(\Lambda^\vee_{\ddelta^\prime}|_{{\calE_{\ddelta^\iota}}} )\chi_\T(W^+_{\d^\iota};\calO_{W^+_{\d^\iota}} )\\ \nonumber&=& m_{\calE_{\d^\prime}}(t^{-1}) m_{\calE_{\d^\iota}}(t)\chi_\T(\Sym(\calA^*))\\ \nonumber&=& \frac{\chi_\T(\det(\calA))}{\chi_\T(\det(T^+_{\calE_{\d^\prime}}))}m_{\calE_{\d^\prime}}(t) m_{\calE_{\d^\iota}}(t)\chi_\T(\Sym(\calA^*)) \\ &=& \frac{\chi_\T(\det(T^{\leq 0}_{\calE_{\d^\prime}}))}{\chi_\T(\det(\C_{1}\otimes\calA^*))}m_{\calE_{\d^\prime}}(t) m_{\calE_{\d^\iota}}(t)\chi_\T(\Sym(\calA^*)), \label{rhs}
\source{very-stable-higgs-bundles-equivariant-multiplicity-mirror-symmetry}
	\eeq
	where in the fourth equation we used \eqref{palindromic}. For the last equation we can argue as follows. As $\iota$ of \eqref{iota} is a $\T$-equivariant isomorphism, it follows that $T^+_{\calE_{\d^\iota}}\cong T^+_{\calE_{\d}}$ and thus $m_{\calE_{\d^\iota}}(t)=m_{\calE_{\d}}(t)$. Using the homogeneity $1$ symplectic form on $T_{\calE_{\d^\prime}} $ we see that $\chi_\T(\C_{-1}\otimes\det(T^+_{\calE_{\d^\prime}}))\chi_\T(\det(T^{\leq  0}_{\calE_{\d^\prime}}))=1$. This implies the last equation. Comparing \eqref{lhs} and \eqref{rhs} implies the first equation in the Theorem.
	
	For the second equation we note that $\dim \calA=\sum_{i=1}^n\dim H^0(C;K^i)= n^2(g-1)+1$ and that a straightforward computation  gives $$\chi_\T(\C_1\otimes \calA^*)=t^{\sum_i (i-1)\dim \calA_i}=t^{\sum_{i=1}^n (i-1)\dim H^0(C;K^i)}=t^{\sum_{i=1}^n (i-1)(2i-1)(g-1)}=t^{(4n+1)(n-1)n(g-1)/6} .$$ The result follows. 
	
	
	
	
	%Recall the type $(1,\dotsm,1)$ fixed point component $F_{{l},\mm}\in \pi_0(\M^{\T})$ from \eqref{type111fix}. We will construct a local universal family $(\bE_{F_{{l},\mm}},\bPhi_{F_{{l},\mm}})$ of Higgs bundles on $C\times F_{{l},\mm}$. Denote by $p:C\times J_{{l}}(C)\to C\times J_0(C)$ the isomorphism sending a degree ${l}$ line bundle $L\in J_{{l}}(C)$ to $L(-{l}c_0)\in J_0(C)$. Then define $$\L_{{l}}:=p^*({\L})\pi_1^{*}(\calO({l}c_0)),$$ which is a universal line bundle on $C\times J_{{l}}(C),$ normalised such that $L_{{l}}|_{c_0\times J}\cong \calO_J$. Let $\mathbb{\Delta}_i\subset C\times C^{[m_i]}$ be the universal divisor, and $$\mathbb{b}_i\in H^0(C\times C^{[m_i]};\calO(\mathbb{\Delta}_i)),$$ the canonical section with $\div(\mathbb{b}_i)=\mathbb{\Delta}_i$. Define now $$\bE_{F_{{l},\mm}}:=\L_{{l}}\oplus \L_{{l}}(\mathbb{\Delta}_1)K^{-1}\oplus \dots\oplus \L_{{l}}(\mathbb{\Delta}_1+\dots+\mathbb{\Delta}_{n-1})K^{-(n-1)}$$ 	and $$\bPhi_{F_{{l},\mm}}:= \left(\begin{array}{ccccc}  0 &0& \dots &0& 0\\ \mathbb{b}_1 & 0 & \dots &0&0\\ 0 & \mathbb{b}_2 & \dots &0&0\\ \vdots & \vdots & \ddots &\vdots &\vdots \\ 0 & 0& \dots &\mathbb{b}_{n-1} & 0
	%\end{array}\right):\bE_{F_{{l},\mm}}\to \bE_{F_{{l},\mm}}\otimes K.$$ Then $(\bE_{F_{{l},\mm}},\bPhi_{F_{{l},\mm}})$ is a family of stable Higgs bundles on $C\times F_{{l},\mm}$ with $\det(\bE_{F_{{l},\mm}})=\L_{{l}}^n(\sum_i (n-i)\mathbb{\Delta}_i)K^{-\binom{n}{2}}$ and thus $\Nm(\bE_{F_{{l},\mm}})=\L_{{l}}^n(\sum_i (n-i)\mathbb{\Delta}_i)$.  
	
	
	%Finally, if we let $\T$ acting on $\L_{{l}}(\Delta_1+\dots \Delta_{i})K^{-i}$ with weight $1-i$, then the $\T$-action on $\bE_{F_{{l},\mm}}$ will induce a weight $-1$ action on $\bPhi_{F_{{l},\mm}}$. Then $(\bE_{F_{{l},\mm}},\bPhi_{F_{{l},\mm}})$ is a family of $\T$-equivariant stable Higgs bundles on $C\times F_{{l},\mm}$ with $\det(\bE_{F_{{l},\mm}})=\L_{{l}}^n(\sum_i (n-i)\mathbb{\Delta}_i)K^{-\binom{n}{2}}$. Forming now $\nm^*(\P_{{L}(-lc_0)})\otimes\bigotimes_{i=1}^{n-1}\bigotimes_{j=1}^{m_i} \Lambda^{n-i}(\bE_{F_{{l},\mm}}|_{\{c_{ij}\}\times F_{{l},\mm} })$ and renormalising the $\T$-action we see that $$\left[\nm^*(\P_{{L}(-lc_0)})\otimes\bigotimes_{i=1}^{n-1}\bigotimes_{j=1}^{m_i} \Lambda^{n-i}(\bE_{c_{ij}})|_{\{\calE_{{L}^\prime;\bb^\prime}\}}\right]^\T=\prod_{i=1}^{n-1} \left[\begin{array}{c} n \\ i\end{array}\right]^{m_i}_t$$
	
	
	
	
	
	
\end{proof}

\begin{remark}
\source{very-stable-higgs-bundles-equivariant-multiplicity-mirror-symmetry}
\notready We proved in Theorem~\ref{fouriermukai} that the relative Fourier-Mukai transform $S(\calL_{\d})=\Lambda_{\d}$ holds over $\calA^\#$ and in Remark~\ref{remarkdual} that  $S(\Lambda_{\d})=\calL_{\d^\iota}[-n^2(g-1)-1]$ also holds over $\calA^\#$. The above theorem  would be a consequence  of a $\T$-equivariant extension of  the mirror symmetry conjecture in the "classical limit"  as in \cite[Conjecture 2.5]{donagi-pantev}. 
\end{remark}

\begin{remark}
\source{very-stable-higgs-bundles-equivariant-multiplicity-mirror-symmetry}
\notready Similar agreement of $\T$-equivariant Euler pairings between certain conjectured mirror pairs of Lagrangian and hyperholomorphic branes of \cite{hitchin2} were observed 
	in \cite[\S 7]{hausel-mellit-pei}, when $n=2$. When $n=2$ we have $\chi_\T(\C_1\otimes \calA^*)=t^{3g-3}$ in Theorem~\ref{euleriso} which we cannot get rid of by just renormalising the $\T$-action on our branes when $3g-3$ is odd. The same phenomenon was observed in \cite[Remark 7.7]{hausel-mellit-pei}. 
	
	In fact, one can also pair certain $n=2$ analogues of the mirror pairs of this paper with the ones in \cite{hausel-mellit-pei} and find the expected agreement. 
	
\end{remark}

\begin{remark}
\source{very-stable-higgs-bundles-equivariant-multiplicity-mirror-symmetry}
\notready Just as in \cite[Corollary 4.11]{hausel-mellit-pei} we get a clearly symmetric expression in $\d$ and $\d^\prime$ if we compute, equivalently as in  \eqref{lhs},
	\beq	\label{tensorsym}\chi_\T(\M;\calL_{\d^\prime}\otimes \Lambda_{\d})=m_{\calE_{\d^\prime}}(t)m_{\calE_{\d}}(t)\chi_\T(\Sym(\calA^*)).\eeq When deriving the agreement of the corresponding Euler forms \cite[Corollary 7.1]{hausel-mellit-pei} we basically argued that it follows from the fact that both of the branes $\calL_i$ and $\Lambda_j$ were self-dual. In our present situation ${\calL_{\d^\prime}}$ is self-dual (up to a shift) for the trivial reason that it is supported on an affine space. On the other hand $\Lambda_{\d}$ is typically not self-dual, rather it satisfies $\Lambda^\vee_{\d}=\iota^*(\Lambda_{\d})=\Lambda_{\d^\iota}$ using the normalised inverse map $\iota$ of Remark~\ref{remarkdual}.  However in \cite{hausel-mellit-pei} both types of our branes $\calL_i$ and $\Lambda_j$ were secretly invariant under $\iota$ as well. In general, what \eqref{skewinv} tells us is that self-dual branes should be mirror to branes invariant under $\vee \circ \iota$. The former is satisfied by $\calL_{\d^\prime}$ and the latter by $\Lambda_{\d}$. 
\source{very-stable-higgs-bundles-equivariant-multiplicity-mirror-symmetry}
\end{remark}

\begin{remark}
\source{very-stable-higgs-bundles-equivariant-multiplicity-mirror-symmetry}
\notready  \label{mirrorate0}
\source{very-stable-higgs-bundles-equivariant-multiplicity-mirror-symmetry}
	It is instructive to take $\delta^\prime=0=(0,0,\dots,0)$ in \eqref{tensorsym}. Then $\calE_{0}=\calE_{\calO_C,0}$ is the canonical uniformising Higgs bundle of  \eqref{uniform}. It has $m_{\calE_0}(t)=1$ and $\calL_{0}\cong \calO_{W^+_0}$ is the structure sheaf of the canonical Hitchin section, while $\Lambda_0\cong \calO_{\M^s}$.   We get $$\chi_\T(\M;\calO_{W_{\delta}^+}\otimes \Lambda_{0})= m_{\calE_{\d}}(t)\chi_\T(\Sym(\calA^*)).$$   On the other hand $$\chi_\T(\M;\calO_{W_0^+}\otimes \Lambda_{\d})=\chi_\T(\calO_{W_0^+};\Lambda_{\d}|_{W_0^+})=\chi_\T(\Lambda_{\d}|_{\calE_0})\chi_\T(\Sym(\calA^*)).$$ Thus $\chi_\T(\M;\calO_{W_{\delta}^+}\otimes \Lambda_{0})=\chi_\T(\M;\calO_{W_0^+}\otimes \Lambda_{\d})$ (or the equivalent Theorem~\ref{euleriso}) implies that \beq \label{mirroriden} \chi_\T(\Lambda_{\d}|_{\calE_0})=m_{\calE_{\d}}(t).\eeq This was observed already in \eqref{eqmultuniv}, but here we derived it as a consequence of mirror symmetry. Indeed such a relationship is expected from the mirror in general. Because  of the restriction property \eqref{fmiden} of the Fourier-Mukai transform to the identity we expect from the fibrewise Fourier-Mukai transform that the mirror of a bounded complex of coherent sheaves $\calF^\bullet$ on $\M$ should restrict to the canonical Hitchin section as $h^*(h_*(\calF^\bullet) )$. The latter was computed for the upward flow of a very stable Higgs bundle in \eqref{hitchinpush} and \eqref{t-character}. Thus \eqref{mirroriden} verifies this expectation from the mirror of $\calL_{\d}$. 
\source{very-stable-higgs-bundles-equivariant-multiplicity-mirror-symmetry}
\end{remark}

\begin{remark}
\source{very-stable-higgs-bundles-equivariant-multiplicity-mirror-symmetry}
\notready We can also generalise these ideas further for very stable upward flows of type not $(1,\dots,1)$. For  $\calE\in \M^{s\T}$  a very stable Higgs bundle let $\calL_\calE:=\calO_{W^+_\calE}$ denote the structure sheaf of its upward flow and  $\Lambda_\calE$   denote its mirror, which we expect to be a $\T$-equivariant vector bundle on $\M^s$ of rank $m_\calE$. If $\calF\in \M^{s\T}$ is another very stable Higgs bundle then denote by $$m_{\calE,\calF}(t):=\chi_\T(\Lambda_\calE|_\calF)\in \Z[t]$$ the $\T$-character  of the fibre $\Lambda_\calE|_\calF$, which we expect to contain only non-positive weights. In particular,  by \eqref{mirroriden} we have $m_{\calE,\calE_0}(t)= m_\calE(t)$ the equivariant multiplicity. As above from mirror symmetry we expect $$\chi_\T(\M;\calL_\calE\otimes \Lambda_\calF)=\chi_\T(\M;\calL_\calF\otimes \Lambda_\calE)$$ which in turn leads to the identity
	\beq\label{mefidentity} m_\calE(t) m_{\calF,\calE}(t)=m_\calF(t) m_{\calE,\calF}(t).\eeq When both $\calE$ and $\calF$ are of type $(1,\dots,1)$ then $m_{\calE,\calF}(t)=m_\calE(t)$ is just the equivariant multiplicity as $\Lambda_\calE|_{\calF}\cong \Lambda_\calE|_{\calE_0}$  and so \eqref{mefidentity} holds. If only $\calF$ is of type $(1,\dots,1)$ we know what $\Lambda_\calF$ should be and how it restricts to any $\calE$ giving a formula for $m_{\calF,\calE}(t)$ only depending on the type of $\calE$. From \eqref{mefidentity} we can deduce $$m_{\calE,\calF}(t)=\frac{m_\calE(t) m_{\calF,\calE}(t)}{m_\calF(t)}.$$ As it is again expected to be a polynomial in $t$ we get more obstructions on $\calE$ being very stable. 
\source{very-stable-higgs-bundles-equivariant-multiplicity-mirror-symmetry}
	
	To illustrate the above take the example of $n=2$ and $d=-2(g-1)$ and let $\calE$ be a very stable Higgs bundle of type $(2)$, i.e. a very stable bundle $E$ together with the zero Higgs field. For some $0\leq i<g$ let $\calF\in F_{l,2i}$ be a very stable type $(1,1)$ Higgs bundle. Then $\Lambda_\calF$ is a tensor product of universal bundles restricted to points of $C$.  By carefully following the normalisation of the $\T$-action on the tensor product of these restricted  universal bundles given by Lemma~\ref{equinormalise} we find that 
	$$m_{\calF,\calE}(t)=\chi_\T(\Lambda_\calF|_\calE)=2^{2i}t^i .$$ By \eqref{mefidentity} we get that
	\beq \label{cotmir} m_{\calE,\calF}(t)=\frac{m_\calE(t) m_{\calF,\calE}(t)}{m_\calF(t)}=\frac{(1+t)^{3g-3} 2^{2i}t^i}{(1+t)^{2i}}=2^{2i}t^i(1+t)^{3g-3-2i}.\eeq Thus even though we do not know (see \ref{cotfib} for more on this) what vector bundle $\Lambda_\calE$ the mirror of the cotangent fibre at $E$  should be we can infer its $\T$-character at type $(1,1)$ fixed points  from \eqref{cotmir}. 
\source{very-stable-higgs-bundles-equivariant-multiplicity-mirror-symmetry}
	
	
\end{remark}
